\documentclass[11pt]{article}

\usepackage[margin=1in]{geometry}
\usepackage{amsmath}
\usepackage{amssymb}
\usepackage{booktabs}
\usepackage[hidelinks]{hyperref}
\usepackage{url}

\title{3D Gaussian Splatting for Temporal AA, Super Resolution, and Motion-Accurate Frame Interpolation}
\author{vkGSplat research note}
\date{Started May 4, 2026}

\begin{document}
\maketitle

\section{Executive Thesis}

The idea is plausible and worth a serious research branch:

\begin{quote}
Use 3D Gaussian Splatting not merely as a renderer, but as an explicit
spatio-temporal reconstruction state for temporal anti-aliasing, temporal
super-resolution, frame generation, and motion-accurate frame interpolation.
The goal is a compute-first counterpart to DLSS 4.5 and later: less of a
black-box image upscaler, more of a renderer whose history buffer is a
persistent, splattable 3D/4D representation.
\end{quote}

As of May 2026, DLSS 4.5 is a real product target: NVIDIA describes it as a
second-generation transformer model for Super Resolution plus Dynamic Multi
Frame Generation and 6X Multi Frame Generation on supported RTX 50-series GPUs
\cite{nvidia2026dlss45,nvidia2026dlssdeveloper}. The research question for
vkGSplat is not whether an independent project can out-train NVIDIA's network.
The sharper question is whether a renderer that owns its scene/path state can
avoid some of the ambiguity that image-space temporal reconstruction must learn
around disocclusion, subpixel geometry, specular motion, and frame interpolation.

Current subjective likelihoods:

\begin{itemize}
  \item Useful vkGSplat temporal reconstruction experiment: \textbf{high}, about 0.8.
  \item 3DGS-assisted TAA beating current CPU reprojection on static scenes: \textbf{medium-high}, about 0.65.
  \item 3DGS-assisted 2x temporal super-resolution prototype: \textbf{medium}, about 0.55.
  \item Motion-accurate 3DGS frame interpolation for rendered/captured clips: \textbf{medium-high}, about 0.65.
  \item 3DGS-based game-style frame generation competitive with DLSS-style MFG: \textbf{low-medium}, about 0.3.
  \item A full ``DLSS 4.5 counterpart'' without a learned image model: \textbf{low}, about 0.2.
\end{itemize}

\section{Why 3DGS Might Matter for TAA/SR}

Modern temporal upscalers operate on a low-resolution, jittered current frame
plus engine-provided auxiliaries: depth, motion vectors, exposure, reactive
masks, and history. DLSS 2 made this temporal-super-sampling pattern neural;
DLSS 4 moved Super Resolution, DLAA, and Ray Reconstruction to transformer
models; DLSS 4.5 upgraded the transformer model and extended frame generation
\cite{nvidia2025dlss4,nvidia2026dlss45}. XeSS is explicitly described by Intel
as AI-based temporal super sampling and anti-aliasing using jittered color,
motion vectors, and depth \cite{intel2026xess}. AMD FSR also evolved from
spatial upscaling to temporal upscaling and frame generation, and the newer FSR
``Redstone'' branding includes upscaling, frame generation, ray regeneration,
and radiance caching \cite{amd2026fsr}.

The common limitation is that most of this reconstruction state is image-space
or screen-derived. That is practical and general, but it means the algorithm
must infer whether a sample is stable history, newly disoccluded content, a
transparent particle, specular reflection, UI, or thin geometry. A 3DGS-based
pipeline can carry explicit object/primitive/Gaussian identity through time.
This gives the reconstruction system a stable world-space address for history.

\section{Relevant Research Landscape}

\subsection{Anti-Aliased 3DGS}

The core 3DGS renderer is not alias-free. At different resolutions or camera
distances, point-sampled Gaussian evaluation can produce blur, jaggies, popping,
and sampling-rate sensitivity. This is directly relevant to temporal SR:
before building a DLSS-like system, the splat renderer itself must understand
pixel footprint.

Multi-Scale 3D Gaussian Splatting keeps Gaussians at multiple scales so distant
or low-resolution rendering can use fewer larger primitives, improving both
quality and speed \cite{yan2024multiscale}. Mip-Splatting is the earlier
must-cite alias-free baseline: it identifies sampling-rate changes, focal
length, and camera distance as sources of 3DGS artifacts, then combines a 3D
smoothing filter with a 2D Mip/box-filter approximation \cite{yu2024mipsplat}.
Analytic-Splatting analytically approximates Gaussian integrals over pixel
windows instead of evaluating at isolated pixel centers
\cite{liang2024analytic}. AAA-Gaussians argues for fuller 3D evaluation
through the pipeline, with adaptive 3D smoothing and stable view-space bounds
to reduce aliasing, projection artifacts, and popping \cite{steiner2025aaa}.
LOD-GS makes filtering strength sampling-rate-sensitive \cite{yang2025lodgs}.
AA-Splat targets feed-forward 3DGS and uses opacity-balanced band-limiting for
arbitrary-resolution rendering \cite{suh2026aasplat}.

These works are mostly spatial or per-frame. The opening for vkGSplat is to
combine their footprint-aware filtering with temporal history.

\subsection{3DGS Super-Resolution}

SplatSuRe studies selective super-resolution for multi-view consistent 3DGS,
motivated by the fact that independently super-resolving input views can
introduce multi-view inconsistency and blurry 3D renders \cite{asthana2026splatsure}.
2D Gaussian Splatting has also been adapted to arbitrary-scale image
super-resolution with feed-forward Gaussian prediction and scale-aware
rasterization \cite{chen2025gassr}. These works support the idea that Gaussian
representations can be used as reconstruction carriers, not only as final
scene geometry.

\subsection{4D and Spacetime Gaussian Splatting}

Frame generation needs a motion model. 4D Gaussian Splatting and related
spacetime Gaussian methods show that Gaussians can carry time and deformation
semantics. 4D-GS uses a deformation field and HexPlane-style structure for
real-time dynamic scenes \cite{wu2024dynamic4dgs}. 4D-Rotor Gaussian Splatting
uses anisotropic XYZT Gaussians and temporal slicing, reporting real-time
rendering rates on RTX hardware \cite{duan2024rotor4dgs}. Spacetime Gaussian
Feature Splatting adds temporal opacity and parametric motion/rotation, with
splatted neural features for dynamic view synthesis \cite{li2023spacetime}.

These papers are not DLSS replacements, but they provide a vocabulary for
Gaussian frame interpolation: interpolate or slice a 4D/temporal Gaussian
state, then splat the intermediate image.

\subsection{Motion-Accurate Frame Interpolation}

This deserves its own category. A normal frame-generation system can be
perceptually smooth while being motion-inaccurate: it may hallucinate an
intermediate image that looks plausible but does not preserve real object
trajectories, depth ordering, or visibility events. For vkGSplat, the stronger
target is \emph{motion-accurate interpolation}: synthesize an intermediate
frame by advancing explicit Gaussian state through a physically or
geometrically meaningful motion model.

RetimeGS is the closest current reference. It frames arbitrary-timestamp
rendering as continuous-time 4DGS reconstruction and identifies poor
inter-frame rendering in prior 4DGS methods as temporal aliasing
\cite{wang2026retimegs}. Its project page emphasizes relatively large
inter-frame motion, coarse flow-guided primitive correspondences,
triple-rendering supervision, spline trajectories, and dynamic stretching as
tools for ghost-free retiming. Its stated limitation is also instructive:
extreme rapid motion can still break optical-flow correspondences.

Other motion-aware Gaussian papers suggest useful mechanisms. Gaussian-Flow
models dynamic Gaussian particles with time/frequency-domain deformation
\cite{lin2023gaussianflow}. GaussianFlow connects Gaussian dynamics and pixel
velocities for 4D content creation \cite{gao2024gaussianflow}. MotionGS uses
explicit motion guidance for deformable 3DGS \cite{zhu2024motiongs}. KF-GS uses
Kalman-filter-style prediction and correction to improve dynamic motion
modeling \cite{liu2026kfgs}. These works point toward a frame interpolation
system whose core data is not only optical flow between two images, but
per-Gaussian position, velocity, covariance, opacity, and visibility confidence.

The broader video-frame-interpolation literature is also useful. AceVFI
classifies VFI methods into center-time frame interpolation and arbitrary-time
frame interpolation, and identifies large motion, occlusion, lighting change,
and non-linear motion as central challenges \cite{kye2026acevfi}. Those are the
same failure modes a 3DGS interpolation system must measure, but 3DGS gives us
a stronger state variable than two images and a flow field.

\subsection{DLSS, FSR, and XeSS as the Bar}

DLSS 3 frame generation used current/prior frames, optical flow, and engine data
such as motion vectors and depth \cite{nvidia2022dlss3}. DLSS 4 introduced
Multi Frame Generation and transformer models for Super Resolution, DLAA, and
Ray Reconstruction \cite{nvidia2025dlss4,nvidia2025dlss4research}. DLSS 4.5
then introduced a second-generation transformer Super Resolution model and
Dynamic Multi Frame Generation / 6X MFG \cite{nvidia2026dlss45}.

This sets a very high bar. A 3DGS counterpart should not begin by trying to
replicate the whole neural stack. It should attack the part DLSS cannot
explicitly own: persistent scene-space reconstruction state.

\section{Candidate Architectures}

\subsection{A: Gaussian History Buffer}

Replace or augment the image-space history buffer with persistent Gaussian
history records:

\[
  H_j =
  (\mu_j, \Sigma_j, q_j, c_j, f_j, v_j, m_j, \tau_j, id_j),
\]

where $\mu_j$ and $\Sigma_j$ are world-space mean/covariance, $q_j$ is
orientation or tangent frame, $c_j$ is accumulated color/radiance, $f_j$ is a
feature vector, $v_j$ is variance, $m_j$ is sample mass, $\tau_j$ is last update
time, and $id_j$ is primitive/Gaussian identity.

For each jittered low-resolution frame:

\begin{enumerate}
  \item unproject visible pixels using depth and camera matrices,
  \item attach primitive ID / material ID / motion vector / confidence,
  \item update nearby Gaussian history entries,
  \item splat the history entries into target-resolution pixels,
  \item blend with current samples using depth, primitive ID, velocity, and
        variance gates.
\end{enumerate}

This is the cleanest first vkGSplat experiment because the repo already has
reprojection contracts in \path{include/vkgsplat/reprojection.h} and
\path{src/core/reprojection.cpp}. The current code reprojects previous color
in screen space and rejects by depth/primitive mismatch. A Gaussian history
buffer keeps the same rejection semantics but moves the address space from
pixels to stable 3D samples.

\subsection{B: Gaussian Temporal Super-Resolution}

Render low-resolution jittered frames. Each frame contributes subpixel samples
to a target-resolution Gaussian reconstruction layer. Instead of resolving
history by clamped pixel neighborhoods, resolve by projected Gaussian footprint:

\[
  C(p) =
  \frac{\sum_j w_j(p, t)\,c_j}
       {\sum_j w_j(p, t) + \epsilon}.
\]

Here $w_j$ should include:

\[
  w_j = w_{\mathrm{footprint}}
        w_{\mathrm{depth}}
        w_{\mathrm{normal}}
        w_{\mathrm{primitive}}
        w_{\mathrm{velocity}}
        w_{\mathrm{variance}}.
\]

The important difference from TAAU is that the accumulated subpixel evidence is
not tied to last frame's integer pixels. It is tied to world-space Gaussian
support and can be re-sampled at the output resolution each frame.

\subsection{C: Gaussian Residual Layer for DLSS-Style Reconstruction}

A hybrid architecture may be strongest:

\[
  \begin{aligned}
  &\text{low-res color/depth/velocity}
  + \text{image history}
  + \text{Gaussian history features}
  \\
  &\rightarrow
  \text{small reconstruction model}.
  \end{aligned}
\]

The model does not need to infer everything from pixels. Gaussian history can
provide stable features:

\begin{itemize}
  \item accumulated radiance,
  \item variance/confidence,
  \item primitive identity consistency,
  \item stable world position,
  \item visibility age,
  \item disocclusion likelihood,
  \item high-resolution subpixel coverage.
\end{itemize}

This is the most plausible route to a DLSS-like result, but it is a later-stage
branch because it needs training data and a small inference runtime.

\subsection{D: Motion-Accurate Gaussian Frame Interpolation}

For interpolation between rendered frames $t_0$ and $t_1$, construct an
intermediate Gaussian state at $t_\alpha$. The simplest model is linear:

\[
  \mu_j(t_\alpha) =
  (1-\alpha)\mu_j(t_0) + \alpha \mu_j(t_1)
\]

with analogous interpolation for rotation, scale, opacity, color/features, and
visibility confidence. Render by splatting the interpolated state.

But a motion-accurate system should move beyond linear interpolation. A better
per-Gaussian state is:

\[
  M_j(t) = (\mu_j(t), \dot{\mu}_j(t), \Sigma_j(t), R_j(t),
            \alpha_j(t), c_j(t), \nu_j(t)),
\]

where $\dot{\mu}_j$ is velocity and $\nu_j$ is visibility confidence. The
trajectory can be constant-velocity, Kalman predicted/corrected, spline-based
as in RetimeGS, or driven by engine object transforms. The renderer should also
produce a forward and backward motion-consistency residual:

\[
  e_j =
  \left\|\Pi_{t_1}(\mu_j(t_1)) -
  \left(\Pi_{t_0}(\mu_j(t_0)) + u_{0\rightarrow1}\right)\right\|,
\]

where $\Pi_t$ projects to the camera at time $t$ and $u_{0\rightarrow1}$ is
engine or optical flow. High residuals mark Gaussians that should not be
trusted for interpolation.

This is the direct 3DGS counterpart to video frame interpolation. It should be
easier for rigid camera motion and object transforms than for particles,
transparency, specular effects, animated shading, or UI. A robust system likely
needs two paths:

\begin{itemize}
  \item geometry-aware Gaussian interpolation for stable scene content,
  \item image-space optical-flow or neural residual generation for everything
        the Gaussian state cannot explain.
\end{itemize}

\subsection{E: Minimal vkGSplat Data Contracts}

The implementation should begin as a reference contract, not a new renderer.
A useful first pair of structs:

\begin{verbatim}
struct GaussianHistorySample {
    float3 position;
    float3 normal;
    float3 color;
    float3 radiance_feature;
    float2 current_to_previous_px;
    float depth;
    float variance;
    uint32_t primitive_id;
    uint32_t material_id;
    uint32_t frame_index;
};

struct GaussianMotionState {
    float3 position_t0;
    float3 position_t1;
    float3 velocity_t0;
    float3 velocity_t1;
    float4 rotation_t0;
    float4 rotation_t1;
    float3 scale_t0;
    float3 scale_t1;
    float opacity_t0;
    float opacity_t1;
    float visibility_confidence;
    float flow_residual_px;
    uint32_t primitive_id;
    uint32_t motion_group_id;
};
\end{verbatim}

\texttt{GaussianHistorySample} is the TAA/SR side: it turns current
reprojection outputs into persistent world-space evidence. \texttt{GaussianMotionState}
is the interpolation side: it stores enough state to synthesize a midpoint and
to reject motion that disagrees with engine or optical flow.

\section{Disocclusion and Ghosting}

The central failure mode of all temporal reconstruction is stale history. The
3DGS advantage is that stale history can be tested in more ways than screen
clamping:

\begin{itemize}
  \item depth agreement after projecting the Gaussian into the current frame,
  \item primitive/Gaussian ID agreement,
  \item normal and tangent-frame agreement,
  \item opacity/transmittance ordering agreement,
  \item current visibility confidence,
  \item age since last observation,
  \item sample variance and neighborhood disagreement.
\end{itemize}

The danger is that 3DGS can make ghosting worse if old Gaussians remain visible
after an object moves or a surface becomes occluded. The cache must support
decay, visibility invalidation, and possibly per-camera frustum evidence.

\section{vkGSplat Experiment Plan}

\subsection{Stage 0: Baseline TAAU Contract}

Extend the current reprojection module into a minimal temporal AA/upscaling
contract:

\begin{itemize}
  \item add Halton or R2 jitter to seed frames,
  \item render low resolution and output high resolution,
  \item reproject history using camera motion map,
  \item add neighborhood clamp or variance clamp,
  \item measure against high-resolution supersampled reference.
\end{itemize}

This is the control. Without it, a Gaussian method has no honest baseline.

\subsection{Stage 1: Static Gaussian History Cache}

Implement CPU-only:

\begin{itemize}
  \item visible pixel $\rightarrow$ world sample,
  \item world samples update Gaussian cache cells,
  \item cache splats to target resolution,
  \item compare against screen-space TAA and current reprojection.
\end{itemize}

Start with static scenes and camera-only motion. Use primitive-ID gates from
the existing vkGSplat reprojection path.

\subsection{Stage 2: 2x Temporal Super-Resolution}

Render at half resolution with jitter and reconstruct full resolution.
Compare:

\begin{itemize}
  \item nearest/bilinear upscale,
  \item screen-space TAAU,
  \item Gaussian history resolve,
  \item Gaussian history plus small spatial filter.
\end{itemize}

Metrics should include PSNR/SSIM/LPIPS/FLIP where available, plus temporal
metrics: shimmer under slow camera pans, ghosting after disocclusion, and edge
stability on thin geometry.

\subsection{Stage 3: Dynamic Gaussian Motion}

Add object transforms and per-Gaussian velocity:

\begin{itemize}
  \item update Gaussian centers/covariances through object motion,
  \item invalidate history on topology/primitive mismatch,
  \item compare against engine motion-vector reprojection,
  \item test on moving rigid objects before deformation.
\end{itemize}

\subsection{Stage 4: Motion-Accurate Gaussian Frame Interpolation}

Generate one intermediate frame between two rendered frames:

\begin{itemize}
  \item interpolate/slice Gaussian state at a chosen $\alpha \in (0,1)$,
  \item start with engine-known camera/object transforms and constant velocity,
  \item add spline/Kalman trajectories for non-linear motion,
  \item validate projected Gaussian motion against engine or optical flow,
  \item splat intermediate view,
  \item generate a disocclusion mask from forward/backward Gaussian visibility,
  \item composite reactive image-space residuals for particles, transparency,
        specular effects, and UI,
  \item compare against optical-flow interpolation and native middle frame.
\end{itemize}

This should stay separate from the TAAU work until Stage 2 is healthy. The
first proof point is not ``more frames'' but whether the generated midpoint
lands moving edges, occlusion boundaries, and object silhouettes in the same
place as a true rendered middle frame.

\subsection{Stage 5: Learned Resolver}

Only after the analytic resolver shows signal:

\begin{itemize}
  \item train a compact model on low-res inputs plus Gaussian history features,
  \item use high-res/high-spp references from vkGSplat or Wicked Engine,
  \item keep inference compatible with CUDA/Metal/Triton-style kernels,
  \item compare with non-Gaussian neural TAAU using the same training data.
\end{itemize}

\section{Key Risks}

\begin{itemize}
  \item The method may become a heavier version of TAA without enough quality
        gain.
  \item Gaussian history can smear across visibility boundaries unless gating
        is strict.
  \item Specular and transparent effects may remain image-space problems.
  \item Dynamic objects need robust lifetime and identity management.
  \item A full DLSS 4.5-class model needs large training data and careful
        latency engineering.
  \item 3DGS visibility ordering can fight temporal stability if splats pop,
        split, merge, or change opacity under optimization.
\end{itemize}

\section{Current Recommendation}

Do pursue this, but sequence it conservatively:

\begin{quote}
First build a strong screen-space TAAU reference. Then add a CPU Gaussian
history cache for static camera-motion scenes. If that beats TAAU on shimmer,
thin geometry, and disocclusion without huge cost, move it to CUDA/Metal and
only then explore learned resolving and frame generation.
\end{quote}

The best vkGSplat framing is not ``we clone DLSS.'' It is:

\begin{quote}
vkGSplat explores explicit 3D/4D reconstruction state as a compute-first
alternative to purely image-space temporal neural rendering.
\end{quote}

\section{Overnight Research TODO}

\begin{itemize}
  \item Read full papers for Multi-Scale 3DGS, Analytic-Splatting,
        AAA-Gaussians, LOD-GS, AA-Splat, SplatSuRe, RetimeGS, and 4DGS variants.
  \item Search for direct ``Gaussian temporal super-resolution'' or
        ``Gaussian TAA'' papers that may not surface under 3DGS keywords.
  \item Compare DLSS 4.5, AMD FSR Redstone, and XeSS 3-style multi-frame
        generation claims using primary sources where possible.
  \item Extract RetimeGS details for continuous-time opacity, spline
        trajectories, flow-guided correspondences, triple rendering, and
        dynamic stretching.
  \item Sketch a \path{GaussianHistoryBuffer} C++ contract compatible with
        current vkGSplat reprojection and denoise frame structs.
  \item Sketch a \path{GaussianMotionState} contract for midpoint synthesis:
        position, velocity, covariance, opacity, visibility, and residual flow.
  \item Identify the minimum benchmark clip needed to expose shimmer,
        disocclusion, thin geometry, and moving-object ghosting.
\end{itemize}

\begin{thebibliography}{99}

\bibitem{nvidia2026dlss45}
NVIDIA.
\newblock DLSS 4.5 Dynamic Multi Frame Generation and Multi Frame Generation
6X Available Now.
\newblock March 31, 2026.
\newblock \url{https://www.nvidia.com/en-us/geforce/news/dlss-4-5-dynamic-multi-frame-generation-6x-mode-released/}

\bibitem{nvidia2026dlssdeveloper}
NVIDIA Developer.
\newblock NVIDIA DLSS.
\newblock Accessed May 4, 2026.
\newblock \url{https://developer.nvidia.com/rtx/dlss}

\bibitem{nvidia2025dlss4}
NVIDIA.
\newblock NVIDIA DLSS 4 Introduces Multi Frame Generation and Enhancements for
All DLSS Technologies.
\newblock January 6, 2025.
\newblock \url{https://www.nvidia.com/en-sg/geforce/news/dlss4-multi-frame-generation-ai-innovations/}

\bibitem{nvidia2025dlss4research}
NVIDIA Research.
\newblock DLSS 4: Transforming Real-Time Graphics with AI.
\newblock Technical report, 2025.
\newblock \url{https://research.nvidia.com/labs/adlr/DLSS4/}

\bibitem{nvidia2022dlss3}
NVIDIA.
\newblock Introducing NVIDIA DLSS 3.
\newblock September 20, 2022.
\newblock \url{https://www.nvidia.com/en-us/geforce/news/dlss3-ai-powered-neural-graphics-innovations/}

\bibitem{amd2026fsr}
AMD.
\newblock AMD FSR Technologies.
\newblock Accessed May 4, 2026.
\newblock \url{https://www.amd.com/en/products/graphics/technologies/fidelityfx/super-resolution.html}

\bibitem{intel2026xess}
Intel.
\newblock Intel XeSS Super Resolution Developer Guide 2.0.
\newblock Accessed May 4, 2026.
\newblock \url{https://www.intel.com/content/www/us/en/developer/articles/technical/xess-sr-developer-guide.html}

\bibitem{yu2024mipsplat}
Zehao Yu, Anpei Chen, Binbin Huang, Torsten Sattler, and Andreas Geiger.
\newblock Mip-Splatting: Alias-Free 3D Gaussian Splatting.
\newblock CVPR, 2024.
\newblock \url{https://openaccess.thecvf.com/content/CVPR2024/html/Yu_Mip-Splatting_Alias-free_3D_Gaussian_Splatting_CVPR_2024_paper.html}

\bibitem{yan2024multiscale}
Zhiwen Yan, Weng Fei Low, Yu Chen, and Gim Hee Lee.
\newblock Multi-Scale 3D Gaussian Splatting for Anti-Aliased Rendering.
\newblock CVPR, 2024.
\newblock \url{https://cvpr.thecvf.com/virtual/2024/poster/30347}

\bibitem{liang2024analytic}
Zhihao Liang, Qi Zhang, Wenbo Hu, Ying Feng, Lei Zhu, and Kui Jia.
\newblock Analytic-Splatting: Anti-Aliased 3D Gaussian Splatting via Analytic
Integration.
\newblock arXiv:2403.11056, 2024.
\newblock \url{https://arxiv.org/abs/2403.11056}

\bibitem{steiner2025aaa}
Michael Steiner, Thomas Koehler, Lukas Radl, Felix Windisch,
Dieter Schmalstieg, and Markus Steinberger.
\newblock AAA-Gaussians: Anti-Aliased and Artifact-Free 3D Gaussian Rendering.
\newblock arXiv preprint, 2025.
\newblock \url{https://tugraz.elsevierpure.com/en/publications/aaa-gaussians-anti-aliased-and-artifact-free-3d-gaussian-renderin/}

\bibitem{yang2025lodgs}
Zhenya Yang, Bingchen Gong, and Kai Chen.
\newblock LOD-GS: Level-of-Detail-Sensitive 3D Gaussian Splatting for Detail
Conserved Anti-Aliasing.
\newblock arXiv:2507.00554, 2025.
\newblock \url{https://papers.cool/arxiv/2507.00554}

\bibitem{suh2026aasplat}
Taewoo Suh, Sungpyo Kim, Jongmin Park, and Munchurl Kim.
\newblock AA-Splat: Anti-Aliased Feed-forward Gaussian Splatting.
\newblock arXiv:2603.29394, 2026.
\newblock \url{https://papers.cool/arxiv/2603.29394}

\bibitem{asthana2026splatsure}
Pranav Asthana, Alex Hanson, Allen Tu, Tom Goldstein, Matthias Zwicker,
and Amitabh Varshney.
\newblock SplatSuRe: Selective Super-Resolution for Multi-view Consistent 3D
Gaussian Splatting.
\newblock CVPR, 2026.
\newblock \url{https://splatsure.github.io/}

\bibitem{chen2025gassr}
Du Chen, Liyi Chen, Zhengqiang Zhang, and Lei Zhang.
\newblock Generalized and Efficient 2D Gaussian Splatting for Arbitrary-scale
Super-Resolution.
\newblock arXiv:2501.06838, 2025.
\newblock \url{https://huggingface.co/papers/2501.06838}

\bibitem{wu2024dynamic4dgs}
Guanjun Wu, Taoran Yi, Jiemin Fang, Lingxi Xie, Xiaopeng Zhang, Wei Wei,
Wenyu Liu, Qi Tian, and Xinggang Wang.
\newblock 4D Gaussian Splatting for Real-Time Dynamic Scene Rendering.
\newblock CVPR, 2024.
\newblock \url{https://guanjunwu.github.io/4dgs/}

\bibitem{duan2024rotor4dgs}
Yuanxing Duan, Fangyin Wei, Qiyu Dai, Yuhang He, Wenzheng Chen, and Baoquan Chen.
\newblock 4D-Rotor Gaussian Splatting: Towards Efficient Novel View Synthesis
for Dynamic Scenes.
\newblock 2024.
\newblock \url{https://research.nvidia.com/publication/2024-02_4d-rotor-gaussian-splatting-towards-efficient-novel-view-synthesis-dynamic}

\bibitem{li2023spacetime}
Zhan Li, Zhang Chen, Zhong Li, and Yi Xu.
\newblock Spacetime Gaussian Feature Splatting for Real-Time Dynamic View
Synthesis.
\newblock arXiv:2312.16812, 2023.
\newblock \url{https://arxiv.org/abs/2312.16812}

\bibitem{wang2026retimegs}
Xuezhen Wang, Li Ma, Yulin Shen, Zeyu Wang, and Pedro V. Sander.
\newblock RetimeGS: Continuous-Time Reconstruction of 4D Gaussian Splatting.
\newblock CVPR, 2026.
\newblock \url{https://william-wang2.github.io/RetimeGS/}

\bibitem{lin2023gaussianflow}
Youtian Lin, Zuozhuo Dai, Siyu Zhu, and Yao Yao.
\newblock Gaussian-Flow: 4D Reconstruction with Dynamic 3D Gaussian Particle.
\newblock arXiv:2312.03431, 2023.
\newblock \url{https://huggingface.co/papers/2312.03431}

\bibitem{gao2024gaussianflow}
Quankai Gao, Qiangeng Xu, Zhe Cao, Ben Mildenhall, Wenchao Ma, Le Chen,
Dahang Tang, and Ulrich Neumann.
\newblock GaussianFlow: Splatting Gaussian Dynamics for 4D Content Creation.
\newblock arXiv:2403.12365, 2024.
\newblock \url{https://huggingface.co/papers/2403.12365}

\bibitem{zhu2024motiongs}
Ruijie Zhu, Yanzhe Liang, Hanzhi Chang, Jiacheng Deng, Jiahao Lu,
Wenfei Yang, Tianzhu Zhang, and Yongdong Zhang.
\newblock MotionGS: Exploring Explicit Motion Guidance for Deformable 3D
Gaussian Splatting.
\newblock arXiv:2410.07707, 2024.
\newblock \url{https://huggingface.co/papers/2410.07707}

\bibitem{liu2026kfgs}
Qingyuan Tang, Yufei Yin, Yanming Zhu, Zhou Yu, Zhenzhong Kuang,
Jiajun Ding, and Jifa He.
\newblock KF-GS: Kalman Filter-Guided Gaussian Splatting for Real-Time
High-Quality Dynamic Scene Reconstruction.
\newblock Journal of Visual Communication and Image Representation, in press,
2026.
\newblock \url{https://www.sciencedirect.com/science/article/pii/S1047320326001100}

\bibitem{kye2026acevfi}
Dahyeon Kye, Changhyun Roh, Sukhun Ko, Chanho Eom, and Jihyong Oh.
\newblock AceVFI: A Comprehensive Survey of Advances in Video Frame
Interpolation.
\newblock IEEE Transactions on Circuits and Systems for Video Technology,
2026.
\newblock \url{https://www.researchgate.net/publication/401799473_AceVFI_A_Comprehensive_Survey_of_Advances_in_Video_Frame_Interpolation}

\end{thebibliography}

\end{document}
